\begin{table}[H]
\centering

\renewcommand{\arraystretch}{1.3} % altura entre filas
\setlength{\tabcolsep}{12pt}      % espacio horizontal entre columnas

\begin{tabular}{ccccc}
\hline
Método & Val. Iniciales & Iteraciones & Raíz $(x)$ & $y(x)$ obtenido \\
\hline

Newton-Raphson & 3 &   - & -       & \text{Error: float() argument must be a str} \\

Newton-Raphson & 1 &   - & -       & \text{Error: float() argument must be a str} \\

Halley         & 3 &   - & -       & \text{Error: float() argument must be a str} \\

Halley         & 1 &   - & -       & \text{Error: float() argument must be a str} \\

Chebyshev      & 3 & 100 & 3.00000 & $1.44225$            \\

Chebyshev      & 1 & 100 & 1.00000 & $1.00000$            \\

Ostrowski      & 3 &   - & -       & \text{Error: float() argument must be a str} \\

Ostrowski      & 1 &   - & -       & \text{Error: float() argument must be a str} \\

Jarratt        & 3 &   - & -       & \text{Error: float() argument must be a str} \\

Jarratt        & 1 &   - & -       & \text{Error: float() argument must be a str} \\

\hline
\end{tabular}

\caption{$y(x) = \sqrt[3]{x}$}
\end{table}


